\subsection{Excluding $C_{74}$}


\begin{proposition}\label{prop:decic74-exclusion}
An elliptic curve over a number field of degree $10$ has no point of order $74$.
\end{proposition}
\begin{proof}
Put $X=X_1(37)$, $J=J_1(37)$ and $\alpha=\langle2\rangle-1$.
Suppose that $E/K$ has a point of order $74$, where $[K:\mathbb Q]=10$.
The point of order $37$ gives a rational effective divisor $D$ of degree
$10$ on $X$, supported away from the cusps, and a $K$-rational
$2$-isogeny. We do not assume initially that the moduli point has degree $10$.
By \cite[Proposition~5.2 and Table~2]{DSPrime}, every diamond-difference
image on $J(\mathbb Q)$ is finite. The exact integral homology certificate
for $X_1(37)$ gives $J(\mathbb Q)[2]=0$: the common binary kernel of
$T_3-1-3\langle3\rangle$ and $T_5-1-5\langle5\rangle$ on the full
rank-$80$ homology is zero. Since the rational torsion reduction kernel
at $2$ is killed by $2$, reduction on $J(\mathbb Q)_{\mathrm{tors}}$ is
injective. We use these degree-independent certificates from the
$X_1(37)$ computation.

The curve has genus $40$ and rational gonality $18$ by
\cite[Theorem~1]{DvH}. Its cusp degrees modulo $2$ are $1^{18},18$;
the eighteen rational cusps lift individually over $\mathbb Q$.
In one global differential basis we compute the actual diamond-$2$
pullback $M$, verify $M^{18}=1$, and put
\[
 U=\operatorname{Im}(M-1),\qquad \dim U=38.
\]
All statements about formal immersion below use the full Jacobian $J$
as target. Canonical base change identifies its cotangent map with
restriction of $U$; no quotient-cotangent exactness at $2$ is required.
The new restriction computations give the following ranks:
\[
\begin{array}{c|r|c}
\text{effective divisor modulo }2&\text{number of tests}&\text{rank}\\\hline
\text{rational cuspidal, degree }10&8\,436\,285&10\\
P_6+\text{rational cuspidal divisor of degree }4&5\,985&10\\
\text{noncuspidal closed point of degree }10&99&10.
\end{array}
\]
Here $P_6$ is one representative of the surviving sextic diamond orbit
in the existing $X_1(37)$ sieve. The other two points are covered by
diamond transport. The cusp jets through order $10$ are transported
by the actual diamond action; their first coefficients identify all
eighteen rational cusps. Two binary algorithms independently check
all multiplicity patterns. Native differential-space calculations check $10c$ and $P_6+4c$.
A fresh reconstruction verifies the actual diamond and Cartier matrices,
all cusp prefix kernels, and the restrictions at $P_6$ and all $99$
decic places.

If $\overline D$ is entirely cuspidal, let $C$ be its global cuspidal
lift. Then $\alpha[D-C]$ is rational torsion reducing to zero, hence
is zero. The first row gives formal immersion of the $\alpha$-Abel
map on $X^{(10)}$ at $\overline C$, so $D=C$, a contradiction.

The complete noncuspidal place counts through degree $10$ are
\[
 (0,0,0,0,0,21,18,0,39,99).
\]
There is consequently at most one noncuspidal component of
$\overline D$; every remaining component is a rational cusp.
Put $A_3=T_3-1-3\langle3\rangle$.
The existing full-source $\alpha A_3$ sieve applies without change:
$A_3$ annihilates rational torsion and kills rational cusps as divisors.
It excludes the septic orbit, one of the two sextic orbits, and two of
the four nonic orbits.

The surviving sextic orbit consists of reductions of three global
sextic divisors above $j=-9317$, as verified by the retained division-
polynomial and Tate-coordinate certificates. If
$\overline D=P_6+C'$ with $\deg C'=4$, choose the corresponding sextic
lift $S$ and the cuspidal lift $C$ of $C'$. Injectivity of torsion
reduction gives $\alpha[D-S-C]=0$. The second row gives formal
immersion at $P_6+C'$, so $D=S+C$, again a contradiction.

For either surviving nonic orbit, its closed point $P_9$ is fixed by
$h=\langle4\rangle$, of order $9$. Write $\overline D=P_9+\overline c$
and lift $\overline c$ to a rational cusp $c$.
The finite diamond-difference image and Manin--Drinfeld imply that
the class of $(h-1)(D-c)$ is rational torsion. It reduces to zero, so
\[
 hD+c\sim D+hc.
\]
Both divisors have degree $11<18$, hence they are equal. Their
noncuspidal parts give $hD=D$. Every geometric orbit of $\langle h\rangle$
on the noncuspidal locus has size divisible by $3$: the projective
automorphism group of an elliptic curve has order at most $3$.
This contradicts $\deg D=10$.

It remains that $\overline D=P_{10}$. In this case the moduli point has
exact degree $10$, and $K$ has a unique prime above $2$, with residue
degree $10$ and ramification index $1$. The curve has good ordinary
reduction there. Indeed, an elliptic curve over $\mathbb F_{1024}$
giving such a point has order $962$ or $1036$, and both traces are odd.
The independent complete curve enumeration gives respectively $5$
and $50$ isomorphism classes, accounting for all $99$ closed level
points.

Let $D'$ be the divisor of the quotient elliptic curve with the image
of the marked $37$-torsion point. In characteristic $2$, the quotient
by a rational subgroup of order $2$ is Frobenius or Verschiebung.
Thus $\overline{D'}=P_{10}$ or $\langle2\rangle P_{10}$.
In the first case $\alpha[D'-D]=0$, and in the second
$\alpha[D'-\langle2\rangle D]=0$, by torsion reduction injectivity.
The third row and formal immersion give respectively
$D'=D$ or $D'=\langle2\rangle D$.
Since $X_1(37)$ is a fine moduli curve, this equality gives a
$K$-rational isogeny of degree $2$ from $E$ to $E^\sigma$ for some
$\sigma\in\operatorname{Aut}(K)$. The unique prime above $2$ is fixed
by $\sigma$. By \cite[Section~4.10]{CCRS}, a geometrically CM elliptic
curve over a decic field cannot have a point of order $37$.
Lemma~\ref{lem:decic-local-two-cycle} now gives a contradiction.
\end{proof}
